% SEGMENT OLP-0499-S01
% part: intuitionistic-logic
% chapter: semantics
% section: introduction

\documentclass[../../../include/open-logic-section]{subfiles}

\begin{document}

\olfileid{int}{sem}{int}

\olsection{அறிமுகம்}

பொருண்மையியல் இல்லாமல் எந்தத் தருக்கமும் நிறைவாக விவரிக்கப்படுவதில்லை;
உள்ளுணர்வுவாதத் தருக்கமும் இதற்கு விதிவிலக்கன்று. செவ்வியல் தருக்கத்திற்கு
!!{valuation}s ஐ அடிப்படையாகக் கொண்ட பொருண்மையியல் நியமமானதாக இருக்க,
உள்ளுணர்வுவாதத் தருக்கத்திற்குப் போட்டியிடும் பல பொருண்மையியல்கள் உள்ளன.
இணைப்பிகளின் பொருள்களுக்கு உள்ளுணர்வுவாத முறையில் ஏற்கத்தக்க விளக்கத்தைத்
தருகின்றன என்ற பொருளில் அவற்றுள் எதுவும் முற்றிலும் நிறைவளிப்பதில்லை.

வாய்ப்புநிலைத் தருக்கங்களுக்கான பொருண்மையியலை ஒத்த, !!{relational model}s ஐ
அடிப்படையாகக் கொண்ட பொருண்மையியலே ஒருவேளை மிகவும் பரவலானது. இப்பொருண்மையியலில்
!!{propositional variable}s உலகங்களுக்கு ஒதுக்கப்படுகின்றன; இவ்வுலகங்கள்
ஓர் அணுகுறவால் தொடர்புபடுத்தப்படுகின்றன. அத்தொடர்பு எப்போதும் ஒரு பகுதி
வரிசை; அதாவது அது தற்சுட்டுப் பண்பும் எதிர்ச்சமச்சீர்ப் பண்பும் கடப்புப்
பண்பும் கொண்டது.

உள்ளுணர்வுசார் முறையில், இவ்வுலகங்களை அறிவு நிலைகளாக அல்லது ``சான்றுநிலைச்
சூழல்களாக'' கருதலாம். நமக்குத் தெரிந்த அளவில் $w'$ ஆனது அறிவின் ஒரு
சாத்தியமான (எதிர்கால) நிலையாக, அதாவது~$w$ இல் அறியப்பட்டவற்றுடன்
ஒத்திசைவானதாக இருந்தால், அப்போதும் அப்போது மட்டுமே, நிலை~$w'$ ஆனது~$w$
இலிருந்து அணுகத்தக்கது. ஒரு கூற்று அறியப்பட்டவுடன், அது மீண்டும்
அறியப்படாததாக முடியாது; அதாவது~$w$ இல் $!A$ அறியப்பட்டு~$Rww'$ எனில்,
$w'$ இலும் $!A$ அறியப்படும். எனவே அணுகுறவைப் பொறுத்து ``அறிவு''
ஓரியக்கமானது.

அறிவுநிலைத் தருக்கத்தில் உள்ளதுபோல் ``$!A$ அறியப்படுகிறது'' என்பதை
``எல்லா அறிவுநிலை மாற்றுகளிலும் மெய்'' என வரையறுத்தால், எல்லா அறிவுநிலை
மாற்றுகளிலும் $!A$,~$!B$ ஆகிய இரண்டும் அறியப்படும்போது~$w$ இல் $!A \land
!B$ அறியப்படுகிறது. ஆனால் அறிவு ஓரியக்கமானதாகவும் $R$ தற்சுட்டுப்
பண்புள்ளதாகவும் இருப்பதால்,~$w$ இல் $!A$, $!B$ ஆகிய இரண்டும்
அறியப்படும்போது, அப்போதும் அப்போது மட்டுமே,~$w$ இல் $!A \land !B$
அறியப்படுகிறது. இதே காரணத்தால், அவற்றுள் குறைந்தது ஒன்று அறியப்படும்போது,
அப்போதும் அப்போது மட்டுமே,~$w$ இல் $!A \lor !B$ அறியப்படுகிறது. எனவே
$\land$, $\lor$ ஆகியவற்றுக்கு இணைப்பிகளின் மெய்மை நிபந்தனைகள் செவ்வியல்
தருக்கத்தில் உள்ளவற்றுடன் ஒன்றுபடுகின்றன.

எனினும், நிபந்தனைவாக்கத்திற்கான மெய்மை நிபந்தனைகள் செவ்வியல்
தருக்கத்திலிருந்து வேறுபடுகின்றன. $Rww'$ ஆகும் எந்த~$w'$ இலும் $!B$ உம்
அறியப்படாமல் $!A$ மட்டும் அறியப்படுவதில்லை எனில், அப்போதும் அப்போது
மட்டுமே,~$w$ இல் $!A \lif !B$ அறியப்படுகிறது. இது~$w$ இல் $!A$
அறியப்படவில்லை அல்லது $!B$ அறியப்படுகிறது என்ற நிபந்தனையைப் போன்றதல்ல.
ஏனெனில்~$w$ இல் $!A$, $!B$ ஆகிய இரண்டையும் நாம் அறியாவிட்டால், $Rww'$
ஆகவும்~$w'$ இல் $!B$ ஐ அறியாமல் $!A$ ஐ மட்டும் அறிந்துகொள்ளும் வகையிலும்
எதிர்கால அறிவுநிலை~$w'$ ஒன்று இருக்கலாம்.

$!A$ ஐ நாம் அறியும் சாத்தியமான எதிர்கால அறிவுநிலை எதுவும் இல்லாதபோது
மட்டுமே $\lnot !A$ ஐ அறிவோம். இங்கு கருத்து இதுதான்: $!A$ அறியத்தக்கதாக
இருந்தால், சாத்தியமான ஏதோவொரு எதிர்கால அறிவுநிலையில்~$!A$ அறியப்படும்.
$\lfalse$ ஐ நாம் அறிய முடியாததால், அந்த எதிர்கால அறிவுநிலையில் $!A$ ஐ
அறிவோம்; ஆனால்~$\lfalse$ ஐ அறிய மாட்டோம்.

இப்பொருள்கோளில் விலக்கப்பட்ட நடுவின் கொள்கை செயலிழக்கிறது. ஏனெனில் சில
$!A$ க்களை நாம் இன்னும் அறியவில்லை; ஆனால் பின்னர் அறிந்துகொள்ளக்கூடும்.
அத்தகைய !!a{formula}~$!A$ க்கு $!A$,~$\lnot !A$ ஆகிய இரண்டும்
அறியப்படாதவை; எனவே $!A \lor \lnot !A$ அறியப்படவில்லை. ஆனால்
எடுத்துக்காட்டாக $\lnot(!A \land \lnot !A)$ என்பதை நாம் அறிவோம். $!A$,
$\lnot !A$ ஆகிய இரண்டையும் அறியும் எதிர்கால நிலை எதுவும் சாத்தியமில்லை;
மேலும் $!A$ அல்லது $\lnot !A$ ஐ நாம் அறிகிறோமா இல்லையா என்பதைச் சாராமல்
இதை அறிவோம்.

உள்ளுணர்வுவாதத் தருக்கத்திற்குக் கிடைக்கும் ஒரே பொருண்மையியல்
!!^{relational model}s அல்ல. இடவியல் பொருண்மையியல் மற்றொன்று: இங்கு
கூற்றுகள் ஓர் இடவியல் வெளியில் திறந்த கணங்களாகப் பொருள்கொள்ளப்படுகின்றன;
இணைப்பிகள் அக்கணங்கள் மீதான செயலிகளாகப் பொருள்கொள்ளப்படுகின்றன
(எடுத்துக்காட்டாக, $\land$ ஆனது வெட்டுக்கு ஒத்துள்ளது).

\end{document}
